%% ============================================================
%% sections/a9-transparency.tex
%% H. R. 9510 Bill v3.0 - APPENDIX E. DEVELOPMENT TRANSPARENCY AND COMMIT SCHEDULE.
%% How Bill v3.0 was built and how a future pass finishes it.
%% Common visuals: Table 9 (version lineage), Table 10 (single-PR commit
%% schedule), Figure 11 (the build-and-handoff process).
%% ============================================================
\clearpage
\phantomsection
\addcontentsline{toc}{section}{Appendix E. Development Transparency and Commit Schedule}
\usccrosshead{APPENDIX E. DEVELOPMENT TRANSPARENCY AND COMMIT SCHEDULE}

\uscprov{1}{\textbf{What was built, and how.} Bill v3.0 (v3.0.0) is the visual
draft of H. R. 9510. It was produced by an autonomous artificial intelligence
coding agent, Claude Code Opus 4.8 (1M context, Max), operating under human
direction in a managed, ephemeral cloud container. It was authored across
sequential commits in a single pull request, one file per commit pushed to GitHub
in real time, with a second-to-last error-fix-and-bundle commit and a final
repository-updates commit. Only kevinkawchak/cancer-automated was edited
\cite{anthropic-claude-code, repo-cancer-automated}.}

\figslot{Table 9.}{Bill Version Lineage}{Bill v1.0 (VVUQ-03 standalone bill), Bill
v2.0 (VVUQ-04 FD\&C Act amendment, v2.0.0 through v2.3.0), and Bill v3.0 (this
VVUQ-05 visual draft, v3.0.0), each with its directory, role, and DOI.}

\uscprov{1}{\textbf{What this version adds.} From Bill v2.0, this version keeps
the full operative text and reorganizes it as a visual perspective: it adds the
figure and table slots and their build instructions, the visual engineering
evidence in Appendix A, the submission deliverables in Appendix B, the
explainability standard in Appendix C, and the commit record here, while confining
the emerging bills and executive actions to Appendix D, exactly as the prior bill
does \cite{kawchak2026vvuq04bill}.}

\figslot{Table 10.}{Single-Pull-Request Commit Schedule}{The ordered commits that
build this bundle: main.tex, the style file, the bibliography, the eight section
files, the README, the error-fix-and-bundle commit with the prompt, output, and
zip, and the final repository-updates commit.}

\uscprov{1}{\textbf{How to finish this measure.} A future Claude Code Opus 4.8 (1M
context) Max pass executes the bracketed drafting instructions in each section
against the named files, rendering every figure and table on a white background as
a text-based diagram and writing each Appendix B deliverable as a separate,
submission-quality Markdown document under papers/VVUQ-05/deliverables. The handoff
is deliberate: this draft names the exact files so the next pass spends its effort
on synthesis, not on searching \cite{vvuq05-figures-bill, vvuq05-next-steps}.}

\figslot{Figure 11.}{The Build-and-Handoff Process}{This v3.0 draft, with its
bracketed instructions and figure and table slots, handed to a future pass that
renders the figures and writes the deliverables, yielding the final visual bill and
the papers/VVUQ-05/deliverables package.}

\begin{draftbox}{Appendix E - the lineage, the commit schedule, and the handoff figure}
\dstep{Table 9 (bill version lineage): build a left-aligned, ragged-right table at the body measure (each width prepended with \rabs) with columns Version, Directory, Role, and DOI. Rows: Bill v1.0, papers/VVUQ-03/final-paper, standalone VVUQ Physical AI Oncology Trial Bill, 10.5281/zenodo.20454870; Bill v2.0, papers/VVUQ-04/final-bill, FD\&C Act amendment (new section 360e-5 and ten conforming amendments), 10.5281/zenodo.20485580; Bill v3.0, papers/VVUQ-05/draft-bill, this visual draft, DOI to be minted. Confirm the directories and DOIs against papers/README.md ("Final Papers and Bills on Zenodo") and papers/VVUQ-04/final-bill/README.md.}
\dstep{Table 10 (commit schedule): build a left-aligned, ragged-right table at the body measure with columns Commit, File or action, and Role. List the fourteen commits of this single pull request in order: (1) main.tex; (2) usctitle.sty; (3) references.bib; (4) sections/s2-findings.tex; (5) sections/s3-amendment.tex; (6) sections/s4-comparative.tex; (7) sections/a5-evidence.tex; (8) sections/a6-deliverables.tex; (9) sections/a7-explainability.tex; (10) sections/a8-research-matrix.tex; (11) sections/a9-transparency.tex; (12) README.md; (13) the error-fix-and-bundle commit adding prompt-draft-bill.md, output-draft-bill.md, and draft-bill-LaTeX.zip; and (14) the final repository-updates commit. Confirm against the actual commit history of this pull request.}
\dstep{Figure 11 (build-and-handoff): render a compact box-and-arrow process in the asciifig environment showing this draft (instructions plus slots) handed to a future pass that renders the figures and writes the deliverables, yielding the final bill and the deliverables package. Reuse the shape of the cover figure in papers/VVUQ-05/draft-bill/main.tex so the two are consistent; keep exact spacing on a white background.}
\dstep{Implementation note: summarize how the measure is implemented as law from papers/VVUQ-04/final-bill/main.tex (Appendix B, the implementation pathway) so the transparency statement and the implementation guidance stay aligned: final regulations within 365 days under section 515D(h), an effective date one year out, a 180-day transition for enrolling investigations, and the readiness gates before enrollment.}
\dstep{Formatting: single hyphens only; the section symbol where a codified reference appears; even RaggedRight spacing with no overflow; keep both tables and the figure on a white background; avoid leaving a single line stranded onto the next page and avoid large empty white bands.}
\end{draftbox}
